% Generated draft: requires mathematical and source review.
\documentclass[11pt]{article}
\usepackage[margin=25mm]{geometry}
\usepackage{fontspec}
\setmainfont{Latin Modern Roman}
\usepackage{amsmath,amssymb,mathtools,unicode-math}
\setmathfont{Latin Modern Math}
\usepackage{enumitem,needspace,xurl,hyperref}
\setlength{\parindent}{0pt}
\setlength{\parskip}{5pt}
\setlength{\emergencystretch}{3em}
\newcommand{\N}{\mathbb{N}}
\newcommand{\Z}{\mathbb{Z}}
\newcommand{\Q}{\mathbb{Q}}
\newcommand{\R}{\mathbb{R}}
\newcommand{\C}{\mathbb{C}}
\newcommand{\F}{\mathbb{F}}
\newcommand{\A}{\mathbb{A}}
\newcommand{\PP}{\mathbb{P}}
\newcommand{\E}{\mathbb{E}}
\newcommand{\eps}{\varepsilon}
\newcommand{\dd}{\,\mathrm{d}}
\begin{document}
\section*{UnsolvedMath ID: 701607}
\textbf{Source problem code:} catalogue entry
\subsection{Definitions and mathematical statement}
Equip \(\Q^3\) with the usual Euclidean distance, expressed by
\[
 d(x,y)^2=\sum_{i=1}^{3}(x_i-y_i)^2.
\]
This is the conventional interpretation of ``isometries'' used below; the supplied record does not explicitly specify the metric. Let \(G=\operatorname{Isom}(\Q^3,d)\) be the group, under composition, of all bijections preserving this distance. For \(g\in G\), define
\[
 \operatorname{Fix}_{\Q^3}(g)=\{x\in\Q^3:g(x)=x\}.
\]
A group is free non-abelian if it is a free group with a free basis containing at least two elements.

The supplied question, proposed by J. Mycielski in 2018 as Kourovka Notebook Problem 19.59, asks whether there exists a free non-abelian subgroup \(H\leq G\) such that every \(h\in H\setminus\{1_H\}\) acts ``without nontrivial fixed points'' in \(\Q^3\). This phrase leaves the following distinction unresolved:
\begin{enumerate}
\item If it means that every nonidentity element has no fixed point at all, the precise question is whether
\[
 \exists H\leq G\quad
 H\text{ is free non-abelian},\qquad
 \forall h\in H\setminus\{1_H\},\quad
 \operatorname{Fix}_{\Q^3}(h)=\varnothing.
\]
\item If ``nontrivial fixed points'' means nonzero fixed vectors, the precise question instead requires
\[
 \exists H\leq G\quad
 H\text{ is free non-abelian},\qquad
 \forall h\in H\setminus\{1_H\},\quad
 \operatorname{Fix}_{\Q^3}(h)\subseteq\{(0,0,0)\}.
\]
\end{enumerate}
The record does not define a trivial fixed point, so these formulations are recorded separately without selecting one as authoritative. Isometries are not assumed to fix the origin or to be linear.

The source labels the rational-space problem open and states that the answer is affirmative when \(\Q\) is replaced by \(\R\), citing Tomkowicz and Wagon. This real-space assertion is supplied background, with the same unresolved fixed-point terminology, and is not independently verified here.
\subsection{Short English statement}
The question asks whether the isometry group of rational Euclidean three-dimensional space contains a free non-abelian subgroup whose nonidentity elements have no nontrivial fixed points. The meaning of ``nontrivial fixed points'' needs confirmation: it may exclude every fixed point or only nonzero fixed points. The supplied source labels the problem open and reports an affirmative real-space analogue.
\subsection{Sources}
\begin{itemize}
\item Local UnsolvedMath registry, permanent ID \(701607\), exact source code \texttt{catalogue entry}, title \emph{Kourovka Notebook Problem 19.59}. The supplied record attributes the problem to J. Mycielski, gives proposal year 2018, and labels it open.
\item \emph{The Kourovka Notebook (Unsolved Problems in Group Theory)}, Problem 19.59. The supplied background identifies the 19th issue (2018) and the 21st issue edition (2026), \texttt{arXiv:1401.0300v48}. The collection description names E. I. Khukhro and V. D. Mazurov as editors.
\item G. Tomkowicz and S. Wagon, \emph{The Banach--Tarski Paradox}, Cambridge Univ. Press, 2016. The supplied statement cites this book for the affirmative analogue over \(\R\); no page or theorem number is supplied.
\end{itemize}
\end{document}
